% SPDX-License-Identifier: CC-BY-4.0
\section{Conclusion}
\label{sec:conclusion}

The absolute position of a lunar surface station is not a quantity that lunar-local measurement can supply on its own: internal ranging leaves a six-dimensional rigid-body defect (Theorem~\ref{thm:rigid}), and that defect closes only through a tie to the Earth/inertial frame. We separated the two ties that can do it. An indirect tie through Earth$\to$satellite ranging reaches the station only when the constellation geometry is rich enough to propagate it, whereas a direct Earth-baseline VLBI tie constrains the station itself regardless of the satellites. The resulting design law is the paper's headline: VLBI \emph{restores} absolute observability when the constellation cannot, and merely \emph{sharpens} it when the constellation already can (Table~\ref{tab:restore}), with the sparse-constellation regime being where VLBI is strictly necessary rather than helpful. The restoration is governed by a sharp threshold of three non-collinear Earth stations, the point at which the residual datum defect reaches zero (\eqref{eq:defect}), and the Cram\'er--Rao bound that follows sets the achievable accuracy along the design curve of Figure~\ref{fig:n_earth}; a Gauss--Newton estimator attains that bound to an efficiency of $1.02$, so the geometry, not the solver, is the limit. These results are exact algebra of the estimation model, computed by an engine that is Validated against independent libraries and the closed-form bounds of \citet{kay1993}, but the lunar network it is applied to remains a representative simulation, so the observability findings themselves stay Modelled; the work is open and fully reproducible. What would turn this analysis into a measurement is a single thing we do not yet have: real lunar VLBI or two-way ranging to an emplaced surface beacon, against which the predicted threshold and bound can be tested.
